\documentclass[11pt]{article}
\usepackage[a4paper,margin=25mm]{geometry}
\usepackage{amsmath,amssymb,amsthm,booktabs,enumitem}
\newtheorem{theorem}{Theorem}
\newtheorem{lemma}[theorem]{Lemma}
\newcommand{\Sym}{\operatorname{Sym}}
\newcommand{\Alt}{\operatorname{Alt}}
\newcommand{\gen}[1]{\langle #1\rangle}
\title{Two endpoint anchors on nine points}
\author{Human Compression project}
\date{}
\begin{document}
\maketitle
\noindent We give a finite boundary calculation for transitive permutation
representations with four marked involutions. Its purpose is to supply the
small case in a separate induction concerning the high-rank representations
studied in \cite{CFL}. The calculation below does not assume that the generated
group is symmetric or alternating, and does not use a pre-existing list of
string C-groups.

Let $\Omega$ be a set of nine points and let $h_0,h_1,h_2,h_3$ be nonidentity
involutions of $\Omega$. Assume the string relations
$[h_i,h_j]=1$ whenever $|i-j|>1$. Write
$H_i=\gen{h_j:j\ne i}$. A transposition of $h_i$ is called \emph{eligible}
if its two points belong to different $H_i$-orbits. An initial anchor is a
point moved by $h_0$ and fixed by the other three marks; a final anchor is
moved by $h_3$ and fixed by the other marks.

\begin{theorem}\label{main}
Suppose that $\gen{h_0,h_1,h_2,h_3}$ is transitive, has distinct initial and
final anchors, and that every mark has at least two eligible transpositions.
Then a relabelling, sending the two anchors to $1$ and $9$, respectively,
puts the ordered tuple in the form
\begin{equation}\label{path}
\begin{split}
h_0&=(1\ 2)(3\ 4),\qquad h_1=(2\ 3)(4\ 5),\\
h_2&=(5\ 6)(7\ 8),\qquad h_3=(6\ 7)(8\ 9).
\end{split}
\end{equation}
This tuple generates the natural $\Alt(9)$ and satisfies the full intersection
property. For $H=\gen{h_1,h_2,h_3}$ its orbits are
\begin{equation}\label{orbits}
\{1\},\quad\{2,3\},\quad\{4,5,6,7,8,9\},
\end{equation}
and the subgroup of $C_{\Sym(9)}(H)$ fixing both anchors is precisely
$\gen{(2\ 3)}$.
\end{theorem}

\begin{proof}[Reduction to eleven endpoint pairs]
Put $a=h_0$ and $d=h_3$. Let the anchors be $b,c$. Because $a$ and $d$
commute, $\{b,b^a\}$ is an orbit on which $d$ acts trivially, and
$\{c,c^d\}$ is an orbit on which $a$ acts trivially. These orbits are disjoint:
otherwise the two orbits coincide, whereas each generator acts nontrivially
on one and trivially on the other. Relabel them as $\{1,2\}$ and $\{8,9\}$,
with $b=1,c=9$. The commuting involutions $a,d$ generate a Klein four-group.
Each orbit on the remaining five points has size one, two, or four.

There are three labelled kinds of orbit of size two: only $a$ moves its
points, only $d$ moves its points, or $a$ and $d$ induce the same swap.
On an orbit of size four the action is regular and is unique up to
relabelling. Let $r_a,r_d,r_b,r_4,f$ be the numbers of these three kinds,
regular orbits, and fixed points. They satisfy
\begin{equation}\label{types}
2(r_a+r_d+r_b)+4r_4+f=5.
\end{equation}
Conversely, every nonnegative solution gives exactly one pair up to
conjugacy fixing $1,2,8,9$. Indeed one may match orbits of each labelled
kind; on a regular orbit a chosen base point determines an equivariant
bijection. Equation~\eqref{types} has the eleven solutions in the table below.
For each solution our representative assigns consecutive labels from
$3,\ldots,7$ first to $a$-only pairs, then $d$-only pairs, then common pairs,
and finally regular orbits. On four consecutive labels $x,y,z,w$ the last
kind uses $a=(x\ y)(z\ w)$, $d=(x\ z)(y\ w)$.
\end{proof}

\noindent It remains to choose $h_1,h_2$. Both fix $1,9$. There are exactly
$7!/(2^j j!(7-2j)!)$ involutions with $j$ swaps on the other seven points;
summing over $j=1,2,3$ gives $231$. For each endpoint representative, let
$B$ be those commuting with $d$ and $C$ those commuting with $a$. We examine
every pair in $B\times C$. This supplies all the remaining string relations.
Repeated marks can be rejected: if two marks coincide, omitting one leaves
the other, so that the omitted mark has no eligible edge.

\begin{center}
\begin{tabular}{rrrrr@{\qquad}rrrr}
\toprule
$r_a$&$r_d$&$r_b$&$r_4$&$f$&$|B|$&$|C|$&transitive&eligible\\
\midrule
0&0&0&0&5&75&75&120&0\\
0&0&0&1&1&11&11&0&0\\
0&0&1&0&3&19&19&0&0\\
0&0&2&0&1&11&11&0&0\\
0&1&0&0&3&19&75&48&0\\
0&1&1&0&1&11&19&8&0\\
0&2&0&0&1&11&75&80&0\\
1&0&0&0&3&75&19&48&0\\
1&0&1&0&1&19&11&8&0\\
1&1&0&0&1&19&19&16&4\\
2&0&0&0&1&75&11&80&0\\
\bottomrule
\end{tabular}
\end{center}
Here the last column imposes at least two eligible edges for every colour,
after transitivity. In all, there are $11\,507$ pairs, $408$ transitive
tuples, and four survivors. These figures are complete finite calculations,
with the following exact certificates and checks.

\begin{enumerate}[label=(\roman*),leftmargin=*]
\item \texttt{enumerate\_boundary.py} generates all permutations of the seven
unanchored points and keeps precisely the nonidentity involutions. For each
of the eleven pairs it tests the two commutation conditions entry by entry.
The connected components of the graph of all marked swaps are exactly the
generated group orbits: each generator preserves these components and each
edge is itself a generator step. The same procedure after omitting each
colour gives exactly its omission orbits. It therefore computes transitivity
and eligibility without a group library or a presupposed generated group.
\item The output gives all four literal surviving tuples and a bijection
from each to \eqref{path}. Each bijection is constructed from $1\mapsto1$ by
propagating all four labelled edges. The certificate checks that every point
gets one value, that the values are distinct, and that every marked edge is
preserved. Both anchors have their specified images. Thus it proves ordered
conjugacy, not merely equality of unlabelled graphs or of group orders.
\item \texttt{endpoint\_control.py} separately enumerates matchings on six
points. After fixing the two distinguished swaps it obtains $76$ choices
for each endpoint and $196$ commuting pairs. The union of the conjugates of
our eleven representatives by all permutations of $\{3,\ldots,7\}$ is
exactly this literal set, with distinct endpoint types disjoint.
\item \texttt{gap\_control.g} independently obtains the $231$ involutions
from the actual pointwise stabilizer of $1,9$ in $\Sym(9)$. It uses actual
subgroup orbits to repeat the full calculation. Its complete survivor set
agrees with the Python list. For \eqref{path}, it then checks equality of the
generated subgroup with the natural $\Alt(9)$, and all $16^2=256$ equalities
\[
 \gen{h_i:i\in I}\cap\gen{h_i:i\in J}=\gen{h_i:i\in I\cap J}
 \qquad(I,J\subseteq\{0,1,2,3\}).
\]
Empty marked subsets are represented by the trivial subgroup. The group
has order $181\,440$, but group equality and intersections are tested
literally; the order alone is not used as an identification.
\end{enumerate}

\begin{proof}[Completion of the proof]
The orbit reduction and the exhaustive middle-pair calculation prove the
classification in Theorem~\ref{main}. The checks in (iv) prove the additional
generation and intersection assertions. The omission graph for
$h_1,h_2,h_3$ in \eqref{path} immediately gives \eqref{orbits}.
Any permutation centralizing $H$ maps $H$-orbits to $H$-orbits of the same
size, so it preserves each of these three orbits. On a transitive orbit an
equivariant permutation fixing a point is the identity: if $x$ is fixed,
then $(x^h)^z=(x^z)^h=x^h$ for every $h\in H$. Consequently a centralizing
permutation fixing $9$ fixes the whole six-point orbit. It also fixes $1$,
and can only swap $2,3$. That swap commutes with each of the three tail marks,
so both possibilities occur. This proves the centralizer assertion.
\end{proof}

\noindent In the notation used for an outer extension, take
$z=1,p=2,q=3,w=4$. The first mark is literally $(z\ p)(q\ w)$ and the only
nonidentity anchored tail centralizer element is $(p\ q)$. No assertion
about the outer extension, about higher degrees, or about the whole source
manuscript is included here. The finite arithmetic is part of this boundary
proof; no earlier small-degree classification is retained as a premise.

\begin{thebibliography}{1}
\bibitem{CFL} P.~J. Cameron, M.~E. Fernandes and D. Leemans,
\emph{The number of string C-groups of high rank},
arXiv:2212.12723v2. This is the context for the separate induction;
its diagrams and tables are not completeness premises here.
\end{thebibliography}
\end{document}
